\chapter{Fees, Rebates and Venue Economics}\label{ch:mx:fees-rebates-and-venue-economics}

The same stock is bid at the same price on two exchanges. One pays the resting order a rebate and has twenty thousand shares queued ahead; the other
charges the resting order a fee and has eight hundred. A trader who wants to be filled should pay to stand in the short line. Exchange fees are a few
hundredths of a cent per share, but the spread of a liquid stock is one cent, and the fee split between the side that rests and the side that takes
changes where orders go, how long queues are and who gets picked off. This chapter prices orders net of fees, shows when the split of a fee matters,
and runs two venues of the exchange simulator side by side, one maker-taker and one inverted, to measure what a rebate buys.

\section{Maker-taker, inverted and flat pricing}

Under maker-taker pricing (One Quant Book 1, chapter~9) the venue pays the resting side a rebate and charges the taking side a fee; an inverted venue
(taker-maker) does the reverse. Both earn the difference.

\begin{definition}[Flat pricing]\label{def:mx:fees-rebates-and-venue-economics:flat}
Under \emph{flat pricing}\index{flat pricing} a venue charges both sides of a trade, or charges them the same fee per share whatever their role,
so that resting and taking cost the same.
\end{definition}

\omterm{def:mx:fees-rebates-and-venue-economics:flat}{Flat pricing} is the norm for futures, where the exchange charges each side a fee per contract, and in many equity markets outside the United
States; maker-taker and inverted pricing are how US equity exchanges compete for the two sides of the order flow within a fee cap on taking displayed
quotes. The two Cboe venues of the box are the chapter's running example; the simulation uses their standard rates.

\begin{dated}{2026-09}{Two fee schedules and the new cap}\label{dat:mx:fees-rebates-and-venue-economics:cboe}
Standard rates for securities at or above USD 1.00, effective September 1, 2026. Cboe BZX (maker-taker): a rebate of USD 0.0016 per share for adding
displayed liquidity, a fee of USD 0.0030 for removing. Cboe BYX (taker-maker): a fee of USD 0.0020 for adding, a rebate of USD 0.0002 for removing.
Both schedules have tiers and fee codes that change these rates for many members. The SEC's 2024 amendments to Rule 610(c) lower the cap on fees for
taking a protected quotation priced at USD 1.00 or more from USD 0.0030 to USD 0.001 per share; exemptive orders of October 31, 2025 and June 11, 2026 deferred compliance to
the first business day of November 2027.
\end{dated}

\section{Fee-adjusted prices and the effective tick}

\begin{definition}[Fee-adjusted price]\label{def:mx:fees-rebates-and-venue-economics:adjusted}
The \emph{fee-adjusted price}\index{fee-adjusted price} of a trade is its price plus the fee per share for a buyer and minus it for a seller, a rebate
counting as a negative fee: what the trade costs or brings once the venue is paid. With $f$ the fee for the trade's role and $\varepsilon=\pm1$ its side,
it is $p+\varepsilon f$.
\end{definition}

A buyer who takes 1\,000 shares at 20.00 pays USD 20.0030 a share on the maker-taker venue and USD 19.9998 on the inverted one: USD 3.20 less on the
inverted venue at the same quoted price. The fee differences create prices between the ticks.

\begin{definition}[Effective tick size]\label{def:mx:fees-rebates-and-venue-economics:efftick}
The \emph{effective tick size}\index{effective tick size} is the smallest difference between the \omterm{def:mx:fees-rebates-and-venue-economics:adjusted}{fee-adjusted prices} at which one side can trade
when the same instrument is quoted on venues with different fees: with take fees $f_1$ and $f_2$ and a tick $\delta_{\mathrm{tick}}$, the fee-adjusted
asks $p+f_1$ and $p+f_2$ interleave with gaps $|f_1-f_2|$ and $\delta_{\mathrm{tick}}-|f_1-f_2|$ (fees reduced modulo the tick).
\end{definition}

With the two Cboe take fees, 0.30 and $-0.02$ cents, the fee-adjusted asks of a one-cent stock sit alternately 0.32 and 0.68 cents apart: a quote on the
inverted venue is a price improvement of 0.32 cents that the one-cent tick forbids on any single venue (\cref{fig:mx:fees-rebates-and-venue-economics:ladder}).
Chao, Yao and Ye (2019) argue that this is why a discrete tick fragments US stock exchanges and disperses their fee structures: venues differentiate by
fees because they cannot by price.

\begin{omfigure}
\begin{tikzpicture}[font=\footnotesize,x=6.5cm]
\draw[->,gray] (-0.12,0) -- (1.5,0);
\node[gray,below] at (0.75,-0.9) {fee-adjusted ask, cents above 20.00};
\foreach \x in {0,0.5,1}{\draw[gray] (\x,0.08) -- (\x,-0.08) node[below] {\x};}
\foreach \x/\l in {-0.02/INV at 20.00,0.98/INV at 20.01}{\fill[omThm] (\x,0) circle (2.5pt); \node[above,omThm] at (\x,0.12) {\l};}
\foreach \x/\l in {0.30/MT at 20.00,1.30/MT at 20.01}{\fill[omDef] (\x,0) circle (2.5pt); \node[above,omDef] at (\x,0.12) {\l};}
\draw[<->] (-0.02,0.75) -- (0.30,0.75) node[midway,above] {0.32};
\draw[<->] (0.30,0.75) -- (0.98,0.75) node[midway,above] {0.68};
\draw[<->] (0.98,0.75) -- (1.30,0.75) node[midway,above] {0.32};
\end{tikzpicture}

\omcaption{Fee-adjusted asks of a stock quoted at 20.00 and 20.01 on a maker-taker venue (take fee 0.30 cents) and an inverted venue (take rebate
0.02 cents), in cents above 20.00. The grid a taker sees is finer than the one-cent tick.}\label{fig:mx:fees-rebates-and-venue-economics:ladder}
\end{omfigure}

A taker compares venues on \omterm{def:mx:fees-rebates-and-venue-economics:adjusted}{fee-adjusted prices} and a router splits a marketable order across them in that order.

\omcode{code/firm/venuefees/firm_venuefees.py}{40}{51}{A marketable order split over the venues' displayed quotes in order of fee-adjusted price, best first.}\label{lst:mx:fees-rebates-and-venue-economics:route}

\section{Is the fee split neutral?}

\begin{definition}[Fee neutrality]\label{def:mx:fees-rebates-and-venue-economics:neutral}
The split of a venue's fee between makers and takers is \emph{neutral}\index{fee neutrality} when moving part of the fee from one side to the other,
the total $F=f^{make}+f^{take}$ held fixed, leaves every trader's \omterm{def:mx:fees-rebates-and-venue-economics:adjusted}{fee-adjusted prices} and the market's outcomes unchanged: the quotes absorb the move.
\end{definition}

\begin{proposition}[Neutrality without a tick]\label{prop:mx:fees-rebates-and-venue-economics:neutral}
Let competitive liquidity providers post the ask $a$ at which a sale breaks even, $a-f^{make}=R$, where $R$ is their reservation price (value, costs
and adverse selection). The taker's \omterm{def:mx:fees-rebates-and-venue-economics:adjusted}{fee-adjusted price} is $a+f^{take}=R+F$, and the maker's net price is $R$: both depend on the total fee only. With a
tick $\delta_{\mathrm{tick}}$ the ask must be $\lceil(R+f^{make})/\delta_{\mathrm{tick}}\rceil\,\delta_{\mathrm{tick}}$, and a change of split
that is not a multiple of the tick changes \omterm{def:mx:fees-rebates-and-venue-economics:adjusted}{fee-adjusted prices}.
\end{proposition}

\begin{proof}
Without a tick, $a=R+f^{make}$ and $a+f^{take}=R+F$; moving $\Delta$ from $f^{take}$ to $f^{make}$ moves $a$ by $+\Delta$ and nothing else. With a
tick, $a$ is the smallest grid price at or above $R+f^{make}$, so the maker keeps a rent $a-f^{make}-R\in[0,\delta_{\mathrm{tick}})$ that depends
on $f^{make}$ modulo the tick.
\end{proof}

Moving 0.20 cents from the take fee (0.30 to 0.10) to the make side (a 0.16 rebate becomes a 0.04 fee) shows the difference. Without a tick, an ask at
10.01 moves to 10.012: the taker still pays 10.013 all-in and the maker still nets 10.0116. On a one-cent grid the ask must stay at 10.01 (the taker now
pays 10.011, the maker loses 0.2 cents) or jump to 10.02 (the taker pays 10.021). When quotes cannot move, the rent goes to queue position instead:
makers compete by joining queues, and the fee split decides how long the queues are. Colliard and Foucault (2012) model how the level and split of fees
change fill rates and spreads; Foucault, Kadan and Kandel (2013) explain maker-taker pricing by the speed at which the two sides react. On the Toronto
Stock Exchange, Malinova and Park (2015) found that posted quotes adjusted to a change in the fee split and liquidity demanders' costs, fees included,
did not change, while posted spreads narrowed and aggressive orders became more frequent.

\section{Tiers and the venue's business}

A venue's revenue from trading is the difference between what takers pay and what makers receive, times its volume; its members' rates depend on
volume tiers (One Quant Book 1, chapter~29), which give the largest liquidity providers the largest rebates and make volume thresholds into cliffs.
Trading fees are only part of the business: market data and connectivity are sold to the same members.

The simulation builds the choice that tiers and fees create. Two engines of \texttt{firm.exchsim} run one stock with a one-cent tick: MT with BZX's
standard rates, INV with BYX's. A small reactive market supplies the flow: an efficient price $P^\ast$ that moves one tick at random times; \omterm{def:mx:the-limit-order-book:orders}{limit
orders} that join the best or rest a few ticks behind and cancel at a constant rate, and at once with probability one half when $P^\ast$ jumps through
them; noise \omterm{def:mx:the-limit-order-book:orders}{market orders}; and informed \omterm{def:mx:the-limit-order-book:orders}{market orders} at a rate proportional to the distance between $P^\ast$ and the mid. Each \omterm{def:mx:the-limit-order-book:orders}{limit order} goes to MT
or INV with probability one half. A fee-aware taker routes by \omterm{def:mx:fees-rebates-and-venue-economics:adjusted}{fee-adjusted price} (\cref{lst:mx:fees-rebates-and-venue-economics:route}), so at equal
prices it takes INV's queue first; a fee-blind one takes the larger displayed queue first, as a router that looks only at size would. Each case runs
three seeds of two hours.

\omcode{code/microstructure/07-fees-rebates-and-venue-economics/python/mx_fees.py}{128}{137}{The takers of the two-venue market: fee-aware routing by fee-adjusted price, fee-blind routing by displayed size.}\label{lst:mx:fees-rebates-and-venue-economics:takers}

When every taker is fee-aware, the inverted venue is the front of a consolidated queue and the maker-taker venue its back. INV takes 77\% of the
volume. Its \omterm{def:mx:the-limit-order-book:orders}{limit orders} receive a fill for 28.2\% of the shares posted, after a median 4.6 seconds, against 8.4\% after 8.8 seconds on MT, where the
queue at the best is longer (910 shares a side on average, against 510). Fills on MT come when INV's queue has been exhausted, which is when the flow is
large or informed: the mid moves 0.63 cents against an MT fill within 30 seconds, 0.13 against an INV fill (standard errors 0.04). The capture is the same
on both, 0.54 cents. Net of adverse selection, an MT fill is worth $-0.09$ cents before the rebate and an INV fill 0.40 cents before the fee.

\begin{omfigure}
\begin{tikzpicture}
\begin{groupplot}[group style={group size=2 by 1,horizontal sep=1.5cm},width=6.6cm,height=5cm,
  tick label style={font=\footnotesize},label style={font=\small},grid=major,grid style={gray!25},table/col sep=comma,
  xtick={0,0.5,1},xlabel={share of fee-aware takers},xmin=-0.1,xmax=1.1]
\nextgroupplot[title={\small filled share of posted shares},ymin=0,ymax=0.32,
  legend columns=2,legend style={font=\footnotesize,at={(1.1,-0.3)},anchor=north,draw=none,fill=none,
  /tikz/every even column/.append style={column sep=8pt}}]
\addplot[omDef,very thick,mark=*] table[x=aware,y=fill_mt]{figdata/microstructure/07-fees-rebates-and-venue-economics/venues.csv};
\addplot[omThm,very thick,mark=square*] table[x=aware,y=fill_inv]{figdata/microstructure/07-fees-rebates-and-venue-economics/venues.csv};
\legend{maker-taker (MT),inverted (INV)}
\nextgroupplot[title={\small 30-second adverse selection (cents)},ymin=0,ymax=0.75]
\addplot[omDef,very thick,mark=*] table[x=aware,y=adverse_mt]{figdata/microstructure/07-fees-rebates-and-venue-economics/venues.csv};
\addplot[omThm,very thick,mark=square*] table[x=aware,y=adverse_inv]{figdata/microstructure/07-fees-rebates-and-venue-economics/venues.csv};
\end{groupplot}
\end{tikzpicture}

\omcaption{Two venues, same flow, half of the \omterm{def:mx:the-limit-order-book:orders}{limit orders} on each. The more takers route by \omterm{def:mx:fees-rebates-and-venue-economics:adjusted}{fee-adjusted price}, the more the inverted venue fills
first and the more the maker-taker venue's fills come from flow that goes on moving the price. Means of three seeds of two hours. Data:
\texttt{mx\_fees.by\_awareness}.}\label{fig:mx:fees-rebates-and-venue-economics:venues}
\end{omfigure}

What does a posted share earn? Its value is the filled share times capture less adverse selection less the make fee. With every taker fee-aware, a
share posted on MT is worth 0.006 cents and one on INV 0.057 (the difference has a standard error of 0.009): the rebate does not pay for the back of the
queue. For a provider to be indifferent, MT would have to pay a rebate of 0.77 cents, more than twice its 0.30-cent take fee: no maker-taker venue can
buy that queue. With half of the takers fee-blind, MT's queue is hit first by those who look at size, and the ranking reverses: 0.052 cents on MT against
0.028 on INV (standard error 0.006), and the indifferent make fee on MT is about zero (0.006 cents), so the 0.16-cent rebate is more than MT needs. With no
fee-aware taker, MT's posted share is worth 0.074 against 0.015. How takers route decides what a rebate buys.

\section{Evidence from fee experiments}

The SEC tried to measure fees as the \omterm{def:mx:the-limit-order-book:tick}{tick size} pilot had measured the tick. It adopted the Transaction Fee Pilot, Rule 610T, on December 19, 2018: 1\,460
randomly selected stocks in two test groups, one with a cap of USD 0.0010 on exchange transaction fees (the cap then in force, set in 2005, was
USD 0.0030) and one in which exchanges could not pay rebates, and every other stock in a control group. Exchanges challenged the rule, and the D.C. Circuit
vacated it on June 16, 2020, before it started. The evidence therefore comes from changes that venues made themselves and from routing data. Malinova
and Park's Toronto study is one; Battalio, Corwin and Jennings (2016) found retail brokers that sent \omterm{def:mx:the-limit-order-book:orders}{limit orders} to the venues paying the largest rebates
and a negative relation between \omterm{def:mx:the-limit-order-book:orders}{limit order} execution quality and the rebate level, the conflict Angel, Harris and Spatt had described. Their finding is
the empirical side of the simulation: an order placed for its rebate stands at the back of the consolidated queue.

\section{Tutorial: the rebate that bought a queue}\label{tut:mx:fees-rebates-and-venue-economics:tut}

\begin{tutorial}
\textbf{Goal.} Price orders net of fees, then measure on two simulated venues what a rebate buys.
\textbf{End state:} \cref{fig:mx:fees-rebates-and-venue-economics:venues} and the numbers of sections 2 to 4.
\begin{tutsteps}
\item \textbf{\omterm{def:mx:fees-rebates-and-venue-economics:adjusted}{Fee-adjusted prices}.} \texttt{firm\_venuefees.fee\_adjusted} and \texttt{effective\_tick} with the box's rates.
\item \textbf{Neutrality.} \texttt{neutral\_ask(ask, take, take\_new, tick)} with and without a tick.
\item \textbf{Two venues.} \texttt{mx\_fees.two\_venues(share\_mt, fees, seconds, seed, flow=Flow(aware=\dots))}: per venue, the filled share, the time to
first fill, capture, adverse selection, value per posted share, queue and volume share.
\item \textbf{Awareness.} \texttt{by\_awareness()} averages three seeds for 0, one half and all takers fee-aware; \texttt{indifference\_make} gives the
make fee that equalises the venues; draw with \texttt{fig\_fees.py}.
\end{tutsteps}
\textbf{What to change next.} Apply the new 0.10-cent cap to MT's take fee and lower its rebate to keep its margin; let each provider choose the venue
with the better expected value given the queues it sees, and watch the shares move.
\end{tutorial}

\section{Build: fees and venue choice}\label{bld:mx:fees-rebates-and-venue-economics:venuefees}

\begin{build}
\bfield{Purpose} Fee-aware prices and venue choice for the smart order router of chapter~18, the transaction-cost analysis of chapter~19 and
Book~11's market makers.
\bfield{Interface} \texttt{fee\_adjusted(price, side, fee)}, \texttt{effective\_tick(tick, fees)}, \texttt{route\_take(quotes, side, qty)},
\texttt{passive\_value(fill, capture, adverse, make)}, \texttt{indifference\_make(\dots)}, \texttt{neutral\_ask(ask, take, take\_new, tick)},
\texttt{tier\_make\_fee(schedule, adav, tcv)} on Book~1's \texttt{firm.feesched}.
\bfield{Rules} Fees per share, positive when paid; a buyer adds the fee, a seller subtracts it; routing sorts quotes by \omterm{def:mx:fees-rebates-and-venue-economics:adjusted}{fee-adjusted price}, displayed
size only.
\bfield{Acceptance tests} \texttt{code/firm/venuefees/tests/}: \omterm{def:mx:fees-rebates-and-venue-economics:adjusted}{fee-adjusted prices} and a split buy and sell by hand; the effective tick for two and
three fees; the indifferent make fee solves its equation; the \omterm{def:mx:fees-rebates-and-venue-economics:neutral}{neutral} ask with and without a tick; a tiered make fee.
\bfield{Stretch} Routing with fill probabilities and queue positions (chapter~6's \texttt{firm.queuevalue}); monthly tier planning with the cliffs of
\texttt{firm.feesched}.
\end{build}

\begin{omsources}
\item J.-E.~Colliard and T.~Foucault, ``Trading fees and efficiency in limit order markets'', \emph{Review of Financial Studies} 25(11), 2012.
\item T.~Foucault, O.~Kadan and E.~Kandel, ``Liquidity cycles and make/take fees in electronic markets'', \emph{Journal of Finance} 68(1), 2013.
\item K.~Malinova and A.~Park, ``Subsidizing liquidity: the impact of make/take fees on market quality'', \emph{Journal of Finance} 70(2), 2015.
\item J.~J.~Angel, L.~E.~Harris and C.~S.~Spatt, ``Equity trading in the 21st century: an update'', \emph{Quarterly Journal of Finance} 5(1), 2015.
\item R.~Battalio, S.~A.~Corwin and R.~Jennings, ``Can brokers have it all? On the relation between make-take fees and limit order execution quality'',
\emph{Journal of Finance} 71(5), 2016.
\item S.~Chao, C.~Yao and M.~Ye, ``Why discrete price fragments U.S. stock exchanges and disperses their fee structures'', \emph{Review of Financial
Studies} 32(3), 2019.
\item United States Court of Appeals for the District of Columbia Circuit, \emph{New York Stock Exchange LLC v. SEC}, No. 19-1042, 2020.
\item Cboe BZX and BYX U.S. Equities fee schedules, September 2026.
\end{omsources}

\section{Exercises}

\begin{exercise}[$\star$]\label{exo:mx:fees-rebates-and-venue-economics:1}
A trader buys 1\,000 shares at 20.00. What is the \omterm{def:mx:fees-rebates-and-venue-economics:adjusted}{fee-adjusted price} with the box's BZX take fee and with BYX's take rebate, and how much does the choice
of venue save?
\end{exercise}

\begin{exercise}[$\star$]\label{exo:mx:fees-rebates-and-venue-economics:2}
What is the effective tick of a one-cent stock traded on two venues whose take fees are 0.30 and $-0.02$ cents?
\end{exercise}

\begin{exercise}[$\star$]\label{exo:mx:fees-rebates-and-venue-economics:3}
A posted share on INV fills 30\% of the time, captures 0.5 cents and loses 0.1 to adverse selection, with a 0.20-cent make fee; on MT it fills 10\% of the
time, captures 0.5 cents and loses 0.6, with a 0.16-cent rebate. What is each worth per posted share?
\end{exercise}

\begin{exercise}[$\star\star$]\label{exo:mx:fees-rebates-and-venue-economics:4}
In exercise~3, what make fee on MT would make a provider indifferent? Can a venue whose take fee is capped at 0.30 cents pay it?
\end{exercise}

\begin{exercise}[$\star\star$]\label{exo:mx:fees-rebates-and-venue-economics:5}
A venue moves 0.20 cents from its 0.30-cent take fee to the make side, where a 0.16-cent rebate becomes a 0.04-cent fee. What happens to an ask at 10.01,
to the taker's \omterm{def:mx:fees-rebates-and-venue-economics:adjusted}{fee-adjusted price} and to the maker's net price, without a tick and with a one-cent tick?
\end{exercise}

\begin{exercise}[$\star\star$]\label{exo:mx:fees-rebates-and-venue-economics:6}
Why does adverse selection on the maker-taker venue rise with the share of fee-aware takers, from 0.31 to 0.63 cents, while the capture stays at 0.54?
\end{exercise}

\begin{exercise}[$\star\star\star$]\label{exo:mx:fees-rebates-and-venue-economics:7}
\emph{Coding.} Run \texttt{two\_venues} with MT's take fee at the new 0.10-cent cap and its rebate cut to 0.06 cents (same margin minus 0.10). With all
takers fee-aware, how do MT's filled share, adverse selection and value per posted share change, and why?
\end{exercise}

\begin{exercise}[$\star\star\star$]\label{exo:mx:fees-rebates-and-venue-economics:8}
\emph{Find the flaw.} ``Our router sends every passive order to the venue that pays the largest rebate, which lowers our clients' costs.''
\end{exercise}

\section{Problem: The Rebate That Bought a Queue}

\begin{problem}[{Weekend problem --- the rebate that bought a queue}]\label{pb:mx:fees-rebates-and-venue-economics:1}
A market maker quotes one stock on a maker-taker and an inverted venue. Decide where to post, and find out what the rebate is paying for.

\medskip\noindent\textbf{Part I --- Prices net of fees.}
\begin{enumerate}
\item Give the box's make and take rates on both venues.
\item Compute the \omterm{def:mx:fees-rebates-and-venue-economics:adjusted}{fee-adjusted price} of a 1\,000-share buy at 20.00 on each.
\item Define the effective tick and compute it for these fees.
\item Why does a discrete tick push venues to compete on fees?
\end{enumerate}

\medskip\noindent\textbf{Part II --- Neutrality.}
\begin{enumerate}[resume]
\item State and prove the neutrality proposition.
\item Rework the move of 0.20 cents from the take fee to the make side with and without a tick.
\item Where does the rent go when quotes cannot move?
\item What did Malinova and Park find on the Toronto Stock Exchange?
\end{enumerate}

\medskip\noindent\textbf{Part III --- Two venues.}
\begin{enumerate}[resume]
\item Describe the simulated market and the two kinds of takers.
\item With all takers fee-aware, give each venue's volume share, filled share and median time to first fill.
\item Give the adverse selection per filled share on each venue and explain the difference.
\item Give the value per posted share on each venue.
\item What rebate would make a provider indifferent, and can MT pay it?
\end{enumerate}

\medskip\noindent\textbf{Part IV --- Routing and the verdict.}
\begin{enumerate}[resume]
\item What changes when half of the takers route by displayed size?
\item And when none are fee-aware?
\item What was the Transaction Fee Pilot, and what happened to it?
\item What did Battalio, Corwin and Jennings find?
\item How does the new access fee cap change the maker-taker venue's room for rebates?
\item State the \emph{named result}: the fill rate and the mark-out of passive orders on the maker-taker and the inverted venue, and the rebate at which a
liquidity provider is indifferent.
\item In one sentence: what does a rebate buy?
\end{enumerate}
\end{problem}

\section{Interview questions}

\begin{interviewq}[$\star$ \iqroles{trader}]\label{iq:mx:fees-rebates-and-venue-economics:1}
What is the difference between a maker-taker and an inverted venue, and where would you send a passive order you need filled quickly?
\end{interviewq}

\begin{interviewq}[$\star\star$ \iqroles{trader, researcher}]\label{iq:mx:fees-rebates-and-venue-economics:2}
Is the split of an exchange fee between makers and takers \omterm{def:mx:fees-rebates-and-venue-economics:neutral}{neutral}? When is it not?
\end{interviewq}

\begin{interviewq}[$\star\star$ \iqroles{researcher}]\label{iq:mx:fees-rebates-and-venue-economics:3}
Passive fills on one venue have worse mark-outs than on another at the same price. Give two explanations and a test that separates them.
\end{interviewq}

\begin{interviewq}[$\star\star$ \iqroles{trader}]\label{iq:mx:fees-rebates-and-venue-economics:4}
Your router ranks venues by \omterm{def:mx:fees-rebates-and-venue-economics:adjusted}{fee-adjusted price}. What does that do to the market makers who post on the most expensive venue to take?
\end{interviewq}

\begin{interviewq}[$\star\star$ \iqroles{researcher}]\label{iq:mx:fees-rebates-and-venue-economics:5}
How would you design an experiment to measure the effect of rebates on market quality?
\end{interviewq}

\begin{interviewq}[$\star\star\star$ \iqroles{trader}]\label{iq:mx:fees-rebates-and-venue-economics:6}
The access fee cap falls from 0.30 to 0.10 cents. What happens to maker-taker rebates, queues and your market-making P\&L?
\end{interviewq}
